\section{ความเป็นมาและความสำคัญของปัญหา}

ภาพและวิดีโอถูกเผยแพร่และส่งต่อได้อย่างรวดเร็วผ่านสื่อสังคมออนไลน์ ระบบกล้องวงจรปิด และพื้นที่จัดเก็บข้อมูลบนคลาวด์ แม้สื่อเหล่านี้มีประโยชน์ต่อการสื่อสารและการวิเคราะห์ข้อมูล แต่ใบหน้าของบุคคลเป็นข้อมูลชีวมิติที่อาจเชื่อมโยงไปยังตัวบุคคลได้โดยตรง การนำสื่อไปใช้ต่อโดยไม่ปกปิดจึงอาจก่อให้เกิดความเสี่ยงด้านความเป็นส่วนตัว

รายงานความก้าวหน้าของภาคการศึกษาก่อนหน้ามุ่งเน้นการออกแบบส่วนติดต่อผู้ใช้และการคัดเลือกโมเดลเบื้องต้น ในรอบปัจจุบันได้พัฒนาระบบให้เป็นกระบวนการทำงานแบบครบลำดับมากขึ้น ตั้งแต่รับไฟล์ภาพหรือวิดีโอ ตรวจจับใบหน้า จัดกลุ่มการปรากฏของบุคคลเดียวกัน เลือกบุคคลที่จะปกปิด สร้างตัวอย่างผลลัพธ์ และส่งออกไฟล์ที่ผ่านการปกปิดแล้ว

ระบบใช้โมเดลตรวจจับแบบ YOLO ซึ่งเป็นแนวคิดการตรวจจับวัตถุแบบขั้นตอนเดียวที่ให้ความสมดุลระหว่างความเร็วและความแม่นยำ ร่วมกับเวกเตอร์ลักษณะใบหน้าจาก ArcFace เพื่อช่วยรวมใบหน้าของบุคคลเดียวกันเป็นรายการเดียวใน Face Map นอกจากการปกปิดอัตโนมัติแล้ว ยังมีเครื่องมือ Manual Control ให้ผู้ใช้วาดกรอบเหนือวัตถุที่ต้องการปกปิด เช่น ป้ายทะเบียนหรือหน้าจอ และเลือกติดตามกรอบดังกล่าวต่อเนื่องในวิดีโอได้

\section{วัตถุประสงค์}

\begin{enumerate}
    \item เพื่อพัฒนาเว็บแอปพลิเคชันสำหรับตรวจจับและปกปิดใบหน้าในไฟล์ภาพและวิดีโอ
    \item เพื่อจัดกลุ่มใบหน้าของบุคคลเดียวกันและให้ผู้ใช้กำหนดว่าจะปกปิดบุคคลใดได้อย่างชัดเจน
    \item เพื่อรองรับการปกปิดข้อมูลด้วยตนเองและการติดตามกรอบสำหรับวัตถุที่อยู่นอกขอบเขตของการตรวจจับอัตโนมัติ
    \item เพื่อให้ผู้ใช้เลือกชนิดและระดับความเข้มของเอฟเฟกต์ปกปิด ตรวจผลลัพธ์ก่อนส่งออก และดาวน์โหลดไฟล์ผลลัพธ์ได้
    \item เพื่อจัดทำแนวทางทดสอบด้านความครอบคลุมของการปกปิด ความแม่นยำ และเวลาในการประมวลผลสำหรับการประเมินในขั้นถัดไป
\end{enumerate}

\section{ขอบเขตและข้อจำกัดของระบบ}

\subsection{ขอบเขต}

ระบบรองรับไฟล์ภาพนามสกุล JPG, JPEG, PNG และ WEBP และไฟล์วิดีโอนามสกุล MP4, MOV, AVI, MKV และ WEBM โดยซอร์สโค้ดกำหนดขนาดไฟล์สูงสุดไว้ 500 MB แต่การบังคับตรวจสอบขนาดก่อนบันทึกยังเป็นส่วนที่ต้องพัฒนาเพิ่มเติม ในด้านการทำงาน ระบบตรวจจับใบหน้าในเฟรมภาพหรือวิดีโอ สร้างรายการบุคคลจากเวกเตอร์ลักษณะ (embedding) และจัดเก็บตำแหน่งที่ตรวจพบไว้ตามเฟรม ผู้ใช้สามารถเลือกปกปิดทุกคน เลือกเฉพาะบางคน หรือปกปิดทุกคนยกเว้นบุคคลที่เลือกได้

การปกปิดรองรับ Gaussian Blur, Pixelation และ Black Box พร้อมค่าความเข้มที่ปรับได้ นอกจากนี้ผู้ใช้ยังวาด ย้าย ปรับขนาด หรือลบกรอบ Manual Blur ได้ กรอบที่เปิดโหมดติดตามจะใช้ตัวติดตาม KCF หรือ CSRT ของ OpenCV ตามความพร้อมของไลบรารี เพื่อพยายามติดตามวัตถุหลังเฟรมเริ่มต้น

\subsection{ข้อจำกัดของโครงงาน}

การตรวจจับอัตโนมัติในปัจจุบันครอบคลุมใบหน้าเป็นหลัก ยังไม่มีแบบจำลองตรวจจับป้ายทะเบียนหรือข้อความส่วนบุคคลแบบอัตโนมัติ จึงต้องใช้ Manual Control กับข้อมูลประเภทดังกล่าว ความแม่นยำของการจัดกลุ่มใบหน้าขึ้นกับคุณภาพของภาพ มุมหน้า แสง การบดบังใบหน้า และค่าเกณฑ์ (threshold) ที่ใช้เปรียบเทียบเวกเตอร์ลักษณะ

การประมวลผลวิดีโอทำแบบเฟรมต่อเฟรม จึงใช้เวลาและทรัพยากรตามความยาวและความละเอียดของวิดีโอ สถานะงานเบื้องหลังถูกเก็บในหน่วยความจำของเซิร์ฟเวอร์ และการอัปโหลดผลลัพธ์ไป Cloudinary ต้องมีการกำหนดข้อมูลรับรองที่ถูกต้องก่อนใช้งานจริง

\section{ประโยชน์ที่คาดว่าจะได้รับ}

\begin{enumerate}
    \item ได้ต้นแบบเครื่องมือปกปิดใบหน้าที่ช่วยลดโอกาสเปิดเผยตัวตนก่อนเผยแพร่ภาพหรือวิดีโอ
    \item ผู้ใช้สามารถควบคุมผลการปกปิดรายบุคคลและปรับรูปแบบการปกปิดให้เหมาะกับบริบทของสื่อ
    \item ได้แนวทางผสานการตรวจจับใบหน้า การรู้จำใบหน้า การติดตามวัตถุ และการประมวลผลวิดีโอเข้ากับเว็บแอปพลิเคชัน
    \item ได้ฐานสำหรับต่อยอดการตรวจจับป้ายทะเบียน ข้อความ หรือข้อมูลส่วนบุคคลประเภทอื่น และสำหรับประเมินประสิทธิภาพเชิงปริมาณต่อไป
\end{enumerate}
